The Vera C. Rubin Observatory began the Legacy Survey of Space and Time (LSST), a ten-year optical survey of the southern sky, on 30 June 2026. This paper presents early Rubin images and discoveries, focusing on the data management system that makes science at this scale possible. Building on  LHC-scale distributed computing experience, Rubin deploys HEP tools, including Rucio and FTS for data movement, CernVM-FS for software distribution and HTCondor for workflow management at its three data facilities: SLAC, an ATLAS Tier-2 site, and the France and UK facilities, both long-standing LHC Tier-1 centres.  We describe the end-to-end Rubin computing ecosystem, present early science results, and discuss lessons from the astronomy--HEP collaboration and open challenges for high-volume, high-rate astronomical surveys.

